\chapter{Mesh-Style Renderers (Legacy Pathways)}\label{ch:mesh-style-renderers}

\section{Overview}

Before the interval shader, rmeshvk rendered tets the classical way: draw
geometry that covers the tet's screen footprint and let the \emph{fragment}
shader intersect the pixel ray with the tet's four planes to find the
entry/exit segment. Four such renderers remain in the tree, selectable in the
viewer via \code{RenderMode}: \code{Regular}, \code{Quad},
\code{MeshShader}, and \code{RayTrace}. They are no longer the default and
are not used by training or the Python forward path, but they are valuable as
cross-checks (the test suite compares renderers against a CPU reference) and
as diagnostics. All of them share the projection + radix-sort front end and
the closed-form volume integral described in the overview article.

Dispatch (\code{crates/rmesh-viewer/src/render.rs}):
\begin{lstlisting}
match render_mode {
  IntervalShader => interval path (see interval-renderer.pdf),
  MeshShader if supported => record_sorted_mesh_forward_pass,
  RayTrace => locate + project + raytrace compute + copy to texture,
  _ /* Regular | Quad */ => record_sorted_forward_pass(use_quad),
}
\end{lstlisting}

\section{Regular: 12 triangles per tet (\code{ForwardPipelines})}

\code{forward_vertex.wgsl} + \code{forward_fragment.wgsl},
\code{record_sorted_forward_pass}. Each visible tet is drawn as an instance
of 12 vertices (4 faces $\times$ 3), with \code{instance_index} indexing the
sorted visible list.

\paragraph{Vertex.} A constant \code{FACE_VERTEX_MAP} picks the face and
corner. The shader computes per-face \emph{inward} normals, the ray--plane
numerators $n \cdot (v_a - o_{\text{cam}})$, and denominators
$n \cdot \text{ray}$, where the ray direction is deliberately left
\emph{unnormalized} as $p_{\text{world}} - o_{\text{cam}}$ so that
perspective-correct interpolation reproduces the exact per-pixel vector.
Back-facing triangles (numerator $\le 0$) are collapsed to zero area ---
only front faces of each tet are rasterized, and each pixel inside the
silhouette is shaded once per front face it covers.

\paragraph{Fragment.} Normalizes the interpolated ray, classifies the four
planes into entering/exiting, and takes
$t_{\min} = \max(t_{\text{enter}})$, $t_{\max} = \min(t_{\text{exit}})$,
$d = \max(t_{\max} - t_{\min}, 0)$; then runs the shared
\code{compute_integral} with the exit color first (webrm convention). MRT
outputs: premultiplied color, an aux target
$(t_{\min}, t_{\max}, \tau, d)$, entry normal, and premultiplied depth. Note
this aux layout differs from the deferred-shading G-buffer --- this path
predates it.

\section{Quad: billboard per tet}

Same record function with \code{use_quad = true}. A compute prepass
(one thread per visible tet, \code{forward_prepass_compute.wgsl}) does the
per-tet work once instead of 12 times per vertex:

\begin{itemize}
\item computes inward face normals and plane numerators;
\item projects the 4 vertices, sorts them by angle around the centroid with a
  5-comparison sorting network to get the convex hull in CCW order;
\item writes a flat 160-byte record per tet: world positions in triangle-strip
  order $[0,1,3,2]$, the four $(\text{normal}, \text{numerator})$ pairs, and
  $(\text{base color}, \text{density})$, $(\nabla c, 0)$.
\end{itemize}

\code{forward_vertex_quad.wgsl} then emits a 4-vertex triangle strip per tet
from the precomputed buffer, and \code{forward_fragment_quad.wgsl} performs
the same 4-plane classification as the Regular fragment shader. It adds a
fallback for behind-camera tets: if the interpolated ray degenerates
(sentinel $\approx 0$), the ray is reconstructed from the pixel coordinates
and camera intrinsics. The quad covers the whole silhouette with one
primitive, trading vertex-shader work for slightly looser coverage.

\section{Mesh-shader forward (\code{MeshForwardPipelines})}

\code{forward_mesh.wgsl} (\code{enable wgpu_mesh_shader}) generates the
per-tet geometry directly in a mesh shader: \code{@workgroup_size(32)} = 8
tets $\times$ 4 threads; thread $k$ of a tet owns vertex $k$ \emph{and} face
$k$. Vertex positions and face normals are exchanged between the 4 threads
with \code{subgroupShuffle} --- groups of 4 never straddle a subgroup
boundary for any subgroup size that is a multiple of 4. Each workgroup emits
32 vertices / 32 primitives; front-facing faces become real triangles,
back-facing ones degenerate. The fragment shader is the unmodified
\code{forward_fragment.wgsl}. Dispatched via \code{indirect_convert.wgsl}
(\code{TETS_PER_GROUP=8}) and \code{draw_mesh_tasks_indirect}. Requires
\code{Features::EXPERIMENTAL_MESH_SHADER}.

\section{Ray-trace compute (\code{rmesh-raytrace})}

\code{RenderMode::RayTrace} is a software ray tracer used as the viewer's
diagnostic mode --- it needs \emph{no sorting} at all, at the price of
$O(\text{depth complexity})$ work per ray.

Flow per frame: the CPU walks the mesh to keep track of which tet contains
the camera (\code{find_containing_tet_walk}); \code{project_compute} runs
only for its SH color evaluation; then \code{raytrace_compute.wgsl}
(\code{@workgroup_size(8,8)}) runs one thread per pixel:

\begin{enumerate}
\item Derive the ray from intrinsics (or from ray buffers when
  \code{ray_mode == 1}).
\item Find the entry tet: if the camera is inside the mesh, start from
  \code{start_tet}; otherwise trace a BVH over the \emph{boundary faces}
  (32-deep stack, M\"oller--Trumbore tests; leaves pack
  \code{tet_id << 2 | face}).
\item \textbf{Adjacency walk}: intersect the ray with the current tet's four
  inward planes (entry $=$ max $t$ over entering faces, exit $=$ min $t$
  over exiting faces), accumulate
  $c \mathrel{+}= c_{\text{premul}} \cdot T_j$ with
  $T_j = e^{\log T}$, plus normal, depth, and entropy moments; step through
  the exit face via \code{tet_neighbors[cur*4 + exit_face]}.
\item Stop on leaving the mesh, on
  $\log T \le -5.54$ (transmittance below $1/255$), or after 512 steps.
\end{enumerate}

The result is written to a storage buffer ($\alpha = 1 - T_{\text{final}}$
plus an aux image with normal, depth, and entropy), then copied to an
\code{Rgba32Float} texture for blitting. \code{locate_compute.wgsl} is a GPU
point-location helper. Because the walk visits tets in exact ray order, this
renderer is the ground truth the rasterizing paths are compared against.

\section{Retained but unwired}

\code{project_compute_16bit.wgsl} (half-precision sort keys, used by
\code{SortMode::Gpu16}) and \code{shadow_ray_gen.wgsl} (a pre-DSM shadow
experiment) are kept in the tree; the latter is superseded by the deep
shadow maps described in Chapter~\ref{ch:shadow-rendering}.

